\section{One-body magic and its decoupling from the sampling cost $|\mathcal S|$}\label{sec:resource}

The value of a sampling method is set by whether the distribution it samples is classically hard, and
for fermions that hardness has a precise algebraic origin. Circuits of free-fermion (matchgate) gates
are classically simulable in polynomial time~\cite{valiant2002}, equivalently the classical
simulability of non-interacting-fermion dynamics~\cite{terhaldivincenzo2002}, with a
Grassmann/covariance-matrix formulation via fermionic linear optics~\cite{bravyi2005flo,jozsamiyake2008}.
The resource that leaves this manifold is fermionic \emph{non-Gaussianity}: every pure non-Gaussian
fermionic state, injected into a matchgate circuit, promotes it to universal quantum
computation~\cite{hebenstreit2019}---the fermionic counterpart of the non-stabilizerness (``magic'')
of universal qubit computation~\cite{bravyikitaev2005,leone2022sre}. A ladder of measures makes this
quantitative: Gaussian rank and extent~\cite{cudbystrelchuk2023}, the non-Gaussian gate count that
bounds phase-sensitive free-fermion simulation~\cite{reardonsmith2024,diaskoenig2024}, and
covariance-matrix monotones computable from two-point functions~\cite{tarabunga2026,collura2026}.

We work with the fermionic AntiFlatness (FAF)~\cite{sierant2026},
\begin{equation}
\mathcal F_k \;=\; 2n_{\rm o}-\tfrac12\,\mathrm{tr}\bigl[(\Gamma^{\top}\Gamma)^k\bigr] \;=\; 2n_{\rm o}-\sum_{j=1}^{2n_{\rm o}}\lambda_j^{2k},
\qquad \Gamma_{mn}=-\tfrac{i}{2}\langle[\gamma_m,\gamma_n]\rangle,
\label{eq:faf}
\end{equation}
where $\Gamma$ is the $4n_{\rm o}\times4n_{\rm o}$ Majorana covariance matrix and $\{\lambda_j\}\in[0,1]$
the singular values of its orthogonal (Youla) normal form---the moduli of its eigenvalue pairs
$\pm i\lambda_j$; here $2n_{\rm o}$ is the number of spin-orbital \emph{modes} (twice the spatial-orbital
count $n_{\rm o}$), distinct from the electron number of the $(N\pm1)$ sectors, and $\mathcal F_k$ ranges in $[0,2n_{\rm o}]$.
$\mathcal F_k=0$ if and only if the state is Gaussian, and $\mathcal F_k$ is invariant
under Gaussian (free-fermion) unitaries by construction~\cite{sierant2026}. The index $k$ is a
hierarchy of correlator orders, efficiently computable and experimentally accessible from two-point
functions~\cite{sierant2026,bittel2025,tarabunga2026}.

\paragraph{One-body collapse.} On the number-conserving states of electronic structure the anomalous
correlators $\langle c_ic_j\rangle$ vanish, so $\Gamma$ is fixed by the one-particle reduced density matrix
$\gamma$ and its singular values are $\lambda_i=|2g_i-1|$, with $g_i\in[0,1]$ the natural
spin-orbital occupations. The $k=1$ member then reduces to a single 1-RDM invariant,
\begin{equation}
\boxed{\;\mathcal F_1 \;=\; \sum_i\bigl[1-(2g_i-1)^2\bigr] \;=\; 4\sum_i g_i(1-g_i)
\;=\; 4\,\mathrm{tr}[\gamma(1-\gamma)]\;,}
\label{eq:collapse}
\end{equation}
where the sum runs over spin orbitals and $g_i\in[0,1]$ are the natural spin-orbital occupations. This
chain is exact on any number-conserving state. For \emph{spin-restricted} natural orbitals
($g_{a\uparrow}=g_{a\downarrow}=n_a/2$) it equals $2N_{\mathrm u}$, with
$N_{\mathrm u}=\sum_a n_a(2-n_a)$ (spatial occupations $n_a\in[0,2]$) the \emph{quadratic}
effectively-unpaired-electron index of Takatsuka--Fueno--Yamaguchi and of
Staroverov--Davidson~\cite{takatsuka1978,staroverov2000} (distinct from the \emph{linear}
$\sum_a\min(n_a,2-n_a)$ index of Head-Gordon~\cite{headgordon2003}); a cleanly singly-occupied spin
orbital contributes $0$ to $\mathcal F_1$ but $2$ to $2N_{\mathrm u}$, so on the five genuinely
spin-polarized (open-shell) members of Table~\ref{tab:molecules} ($\mathrm{O_2},\mathrm{NO},\mathrm{CN},
\mathrm{OH},\mathrm{BeH}$) the two differ and we report the exact
$\mathcal F_1=4\,\mathrm{tr}[\gamma(1-\gamma)]$ directly. The invariant $\mathrm{tr}[\gamma(1-\gamma)]$ is
the Gigena--Rossignoli one-body-entanglement linear entropy~\cite{gigena2020}. We verify
Eq.~\eqref{eq:collapse}---and, on the closed-shell members, its equality to $2N_{\mathrm u}$---to machine
precision across the dissociation of $\mathrm{N_2}$ and the full suite of Table~\ref{tab:molecules}. Algebraically Eq.~\eqref{eq:collapse} is an immediate
number-conserving specialization of the antiflatness of Ref.~\cite{sierant2026}, and its ensemble
form---$\mathcal F_1$ proportional to the one-body purity deficit---was independently noted for
interacting-integrable eigenstates in Ref.~\cite{swietek2026}; our point is not the algebra but the
exact, state-by-state \emph{dictionary} it makes explicit. Three literatures---fermionic magic,
effectively-unpaired electrons, and one-body entanglement---thus name one quantity, and, being a
function of the occupation spectrum alone, $\mathcal F_1$ inherits the orbital-rotation invariance of
the eigenvalues of $\gamma$. (Outside
number conservation the collapse fails: a paired Gaussian state hides fractional occupations in
anomalous correlators and keeps $\mathcal F_1=0$.)

\paragraph{Decoupling from the cost.} That invariance is exactly what severs $\mathcal F_1$ from the
sampling cost. The cost of a sample-based method is the determinant support $|\mathcal S|_\varepsilon$:
the number of fixed-basis Slater determinants that carry $1-\varepsilon^2$ of the weight (concretely,
$99\%$ of the weight for the molecular suite and $\varepsilon=0.05$ for the lattice families). Throughout,
$|\mathcal S|$ without a subscript denotes this thresholded count; where the \emph{raw} support---the
determinants of nonzero amplitude---is meant we say so explicitly, because the two do not obey the same
bounds. Unlike
$\mathcal F_1$, $|\mathcal S|$ is \emph{basis dependent}---a single Slater determinant has
$|\mathcal S|=1$ in its own natural orbitals and exponentially many terms in a generic basis. An
orbital rotation is a free (Gaussian) operation that leaves the occupation spectrum, hence
$\mathcal F_1$, exactly fixed while changing $|\mathcal S|$---from $O(1)$ in the natural-orbital basis to
exponentially many determinants in a generic basis. A quantity invariant under operations that change
$|\mathcal S|$ at fixed value cannot bound or predict it: one-body magic and the sampling cost are
decoupled by symmetry.

\begin{figure}[t]\centering
% ===== fig:decoupling (REBUILT 2026-09-18) -- native pgfplots, no external .dat, no hand-typed numbers.
% Emitted by p2_figs/make_figs.py from data/cost_vs_ent.json, data/n19_suite.json,
% data/resource_master.json and p2_calc/stats_resource.json; all statistics re-verified at emit time.
\definecolor{dcCov}{HTML}{E8543A}\definecolor{dcMrf}{HTML}{10A87B}\definecolor{dcIon}{HTML}{2E6FE0}
\definecolor{dcTrend}{HTML}{9A9A9A}\definecolor{dcPool}{HTML}{6E6E6E}\definecolor{inkPrim}{HTML}{1B1B1B}
\definecolor{scAh}{HTML}{D55E00}\definecolor{scAd}{HTML}{F0A868}\definecolor{scBh}{HTML}{00755E}\definecolor{scBd}{HTML}{5FC2A5}\definecolor{scCh}{HTML}{1B4F9C}\definecolor{scCd}{HTML}{79AEE8}
\begin{tikzpicture}
\begin{groupplot}[group style={group size=2 by 2, horizontal sep=1.8cm, vertical sep=2.05cm},
  width=7.45cm, height=5.6cm, paperaxis,
  title style={at={(0,1)},anchor=south west,font=\small\bfseries,yshift=2pt},
  label style={font=\normalsize}, tick label style={font=\footnotesize},
  ymajorgrids, grid style={draw=black!8}]
% ---- (a) molecules: F1 and |S| rise together (the regime where F1 looks predictive) ----
\nextgroupplot[title={(a)~molecules: they rise together},
  xlabel={$\mathcal F_1=4\,\mathrm{tr}[\gamma(1{-}\gamma)]$}, ylabel={determinant support $|\mathcal S|$},
  ymode=log, xmin=-0.9, xmax=12.29, ymin=0.62, ymax=460, enlarge x limits=false]
\addplot[dcTrend,line width=1pt,dashed,domain=0.000:11.377,samples=60,forget plot] {2^(0.4736*x+0.8502)};
\addplot[only marks,mark=*,mark size=1.9pt,draw=white,line width=0.4pt,fill=dcCov] coordinates {(0.0380,1) (2.5114,2) (0.0008,1) (3.4603,4) (0.0676,1) (3.9849,2) (0.6882,8) (8.1924,46) (0.8474,7) (11.3765,80) (0.0392,1) (3.7542,6) (0.0068,1) (6.9358,13) (0.0803,1) (10.2890,68) (0.2515,2) (5.1131,11) (0.5214,9) (7.7351,91)};
\addplot[only marks,mark=*,mark size=1.9pt,draw=white,line width=0.4pt,fill=dcMrf] coordinates {(4.7428,13) (8.0009,6) (0.0127,1) (3.7026,5) (1.6728,11) (6.0001,4) (1.3066,27) (5.9378,9) (0.7900,5) (7.4487,25)};
\addplot[only marks,mark=*,mark size=1.9pt,draw=white,line width=0.4pt,fill=dcIon] coordinates {(0.0021,1) (0.0750,2) (0.0000,1) (0.0007,1) (0.5181,5) (0.5731,3) (0.0353,1) (0.3864,2)};
\node[inkPrim,font=\scriptsize,anchor=north west,align=left] at (rel axis cs:0.02,0.97) {Pearson $r=0.80$ on $\log_2|\mathcal S|$\\Spearman $\rho=0.862$ ($n=38$)};
% ---- (b) the six (L,N) strata against F1 ----
\nextgroupplot[title={(b)~Hubbard: $\mathcal F_1$ runs backwards},
  xlabel={one-body magic $\mathcal F_1$}, ylabel={determinant support $|\mathcal S|$},
  ymode=log, xmin=-0.9, xmax=20.2, ymin=40, ymax=900000, enlarge x limits=false]
\addplot[scAh,mark=*,mark size=2.1pt,line width=1.0pt,mark options={fill=scAh,draw=scAh}] coordinates {(0.6258,330) (2.2316,312) (5.9132,260) (9.6047,147) (11.2997,68)};
\addplot[scAd,mark=o,mark size=2.1pt,line width=1.0pt,mark options={fill=scAd,draw=scAd}] coordinates {(0.3396,200) (1.1047,196) (2.8100,188) (4.9945,168) (6.5809,143)};
\addplot[scBh,mark=square*,mark size=2.0pt,line width=1.0pt,mark options={fill=scBh,draw=scBh}] coordinates {(0.8020,3395) (2.9184,3120) (7.8807,2176) (12.8163,880) (15.0697,327)};
\addplot[scBd,mark=square,mark size=2.0pt,line width=1.0pt,mark options={fill=scBd,draw=scBd}] coordinates {(0.4855,2342) (1.6348,2272) (4.3683,2041) (7.9600,1600) (10.4394,1091)};
\addplot[scCh,mark=triangle*,mark size=2.6pt,line width=1.0pt,mark options={fill=scCh,draw=scCh}] coordinates {(0.9724,35678) (3.6002,30975) (9.8568,17732) (16.0306,5084) (18.8404,1476)};
\addplot[scCd,mark=triangle,mark size=2.6pt,line width=1.0pt,mark options={fill=scCd,draw=scCd}] coordinates {(0.6329,26503) (2.1803,25235) (5.9969,21203) (11.0065,13832) (14.3023,7527)};
\node[inkPrim,font=\scriptsize,anchor=north west,align=left] at (rel axis cs:0.02,0.97) {$U=1\to16$ along each curve\\within each sweep: $\rho=-1.000$\\pooled ($n=30$): $\rho=-0.109$ (withdrawn)};
% ---- (c) the same six strata against chi ----
\nextgroupplot[title={(c)~Hubbard: $\chi$ runs forwards},
  xlabel={minimal bond dimension $\chi$},
  ylabel={determinant support $|\mathcal S|$},
  ymode=log, xmin=2.6, xmax=18.4, ymin=40, ymax=900000, enlarge x limits=false,
  legend style={at={(1.16,-0.46)},anchor=north,legend columns=6,draw=black!25,font=\footnotesize},
  legend cell align=left]
\addplot[scAh,mark=*,mark size=2.1pt,line width=1.0pt,mark options={fill=scAh,draw=scAh}] coordinates {(12,330) (10,312) (6,260) (5,147) (4,68)};
\addlegendentry{$(6{,}6)$}
\addplot[scAd,mark=o,mark size=2.1pt,line width=1.0pt,mark options={fill=scAd,draw=scAd}] coordinates {(11,200) (10,196) (9,188) (9,168) (9,143)};
\addlegendentry{$(6{,}4)$}
\addplot[scBh,mark=square*,mark size=2.0pt,line width=1.0pt,mark options={fill=scBh,draw=scBh}] coordinates {(11,3395) (11,3120) (8,2176) (8,880) (7,327)};
\addlegendentry{$(8{,}8)$}
\addplot[scBd,mark=square,mark size=2.0pt,line width=1.0pt,mark options={fill=scBd,draw=scBd}] coordinates {(16,2342) (16,2272) (15,2041) (14,1600) (13,1091)};
\addlegendentry{$(8{,}6)$}
\addplot[scCh,mark=triangle*,mark size=2.6pt,line width=1.0pt,mark options={fill=scCh,draw=scCh}] coordinates {(17,35678) (14,30975) (10,17732) (6,5084) (6,1476)};
\addlegendentry{$(10{,}10)$}
\addplot[scCd,mark=triangle,mark size=2.6pt,line width=1.0pt,mark options={fill=scCd,draw=scCd}] coordinates {(13,26503) (13,25235) (13,21203) (11,13832) (9,7527)};
\addlegendentry{$(10{,}8)$}
\node[inkPrim,font=\scriptsize,anchor=north west,align=left] at (rel axis cs:0.02,0.97) {$\chi$ falls as $U$ rises\\within each sweep: $\rho=0.894$ to $1.000$\\pooled ($n=30$): $\rho=+0.569$ (withdrawn)};
% ---- (d) Simpson panel ----
\nextgroupplot[title={(d)~stratified vs pooled},
  xlabel={stratum $(L,N)$}, ylabel={Spearman $\rho$ vs $|\mathcal S|$},
  xmin=0.45, xmax=6.55, ymin=-1.42, ymax=1.95, ytick={-1,-0.5,0,0.5,1},
  xtick={1,2,3,4,5,6}, xticklabels={$(6{,}6)$,$(6{,}4)$,$(8{,}8)$,$(8{,}6)$,$(10{,}10)$,$(10{,}8)$},
  x tick label style={font=\scriptsize,rotate=30,anchor=east}, ymajorgrids=false]
\addplot[dcPool,densely dashed,line width=0.9pt,forget plot] coordinates {(0.45,0.5692) (6.55,0.5692)};
\addplot[dcPool,densely dashed,line width=0.9pt,forget plot] coordinates {(0.45,-0.1088) (6.55,-0.1088)};
\addplot[only marks,mark=*,mark size=2.4pt,scAh,mark options={fill=scAh,draw=white,line width=0.4pt},forget plot] coordinates {(1,1.0000)};
\addplot[only marks,mark=square*,mark size=2.3pt,scAh,mark options={fill=scAh,draw=white,line width=0.4pt},forget plot] coordinates {(1,-1.0000)};
\addplot[only marks,mark=*,mark size=2.4pt,scAd,mark options={fill=scAd,draw=white,line width=0.4pt},forget plot] coordinates {(2,0.8944)};
\addplot[only marks,mark=square*,mark size=2.3pt,scAd,mark options={fill=scAd,draw=white,line width=0.4pt},forget plot] coordinates {(2,-1.0000)};
\addplot[only marks,mark=*,mark size=2.4pt,scBh,mark options={fill=scBh,draw=white,line width=0.4pt},forget plot] coordinates {(3,0.9487)};
\addplot[only marks,mark=square*,mark size=2.3pt,scBh,mark options={fill=scBh,draw=white,line width=0.4pt},forget plot] coordinates {(3,-1.0000)};
\addplot[only marks,mark=*,mark size=2.4pt,scBd,mark options={fill=scBd,draw=white,line width=0.4pt},forget plot] coordinates {(4,0.9747)};
\addplot[only marks,mark=square*,mark size=2.3pt,scBd,mark options={fill=scBd,draw=white,line width=0.4pt},forget plot] coordinates {(4,-1.0000)};
\addplot[only marks,mark=*,mark size=2.4pt,scCh,mark options={fill=scCh,draw=white,line width=0.4pt},forget plot] coordinates {(5,0.9747)};
\addplot[only marks,mark=square*,mark size=2.3pt,scCh,mark options={fill=scCh,draw=white,line width=0.4pt},forget plot] coordinates {(5,-1.0000)};
\addplot[only marks,mark=*,mark size=2.4pt,scCd,mark options={fill=scCd,draw=white,line width=0.4pt},forget plot] coordinates {(6,0.8944)};
\addplot[only marks,mark=square*,mark size=2.3pt,scCd,mark options={fill=scCd,draw=white,line width=0.4pt},forget plot] coordinates {(6,-1.0000)};
\node[inkPrim,font=\scriptsize,anchor=north west] at (rel axis cs:0.02,0.995) {$\bullet$\; $\rho(\chi,|\mathcal S|)$ per stratum};
% D12: this key used to sit at the bottom-left, corner-to-corner with the (6,6) datum.
\node[inkPrim,font=\scriptsize,anchor=north west] at (rel axis cs:0.02,0.90) {$\blacksquare$\; $\rho(\mathcal F_1,|\mathcal S|)$ per stratum};
\node[dcPool,font=\scriptsize,anchor=north west] at (rel axis cs:0.035,0.578) {pooled ($n=30$): $\rho=+0.569$ (withdrawn)};
\node[dcPool,font=\scriptsize,anchor=north west] at (rel axis cs:0.035,0.377) {pooled ($n=30$): $\rho=-0.109$ (withdrawn)};
\end{groupplot}
\end{tikzpicture}

\caption{\label{fig:decoupling}\textbf{One-body magic $\mathcal F_1$ is decoupled from the sample-based
determinant support $|\mathcal S|$.} \textbf{(a)}~Nineteen FCI-verified molecules at equilibrium and
dissociation: in the static-correlation regime $\mathcal F_1$ and $|\mathcal S|$ rise together
(Pearson $r=0.80$ on $\log_2|\mathcal S|$), yet $\mathcal F_1$ approaches its maximum while
$|\mathcal S|$ stays a modest fraction of the active-space sector. \textbf{(b)}~Hubbard chains at fixed
filling (three representative sweeps; the full $30$-point set is Fig.~\ref{fig:master}): increasing on-site repulsion $U$ raises $\mathcal F_1$ but \emph{lowers} the site-basis
support $|\mathcal S|$ as the ground state localizes, inverting the sign of the correlation.
\textbf{(c)}~Free fermions ($U=0$): a single Slater determinant with $\mathcal F_1=0$, whose site-basis
support nonetheless grows as $35,336,3496,37361$ for chain length $L=4,6,8,10$. The
$\mathcal F_1$--$|\mathcal S|$ relationship is not sign stable---the signature of comparing an
orbital invariant to a basis variant.}
\end{figure}

The decoupling is two-sided, and its sign is not even stable (Fig.~\ref{fig:decoupling}). Across the
nineteen FCI-verified molecules of Table~\ref{tab:molecules}, $\mathcal F_1$ and $|\mathcal S|$ rise
\emph{together} in the static-correlation regime (Pearson $r=0.80$ on $\log_2|\mathcal S|$): bond
dissociation drives fractional occupations and multireference weight in step. But a free-fermion
ground state---a single Slater determinant, $\mathcal F_1=0$---has $|\mathcal S|$ that grows
exponentially in the site basis ($|\mathcal S|=35,\,336,\,3496,\,37361$ for chain length
$L=4,6,8,10$), and along a Hubbard chain at fixed filling, increasing on-site repulsion \emph{raises}
$\mathcal F_1$ while \emph{lowering} $|\mathcal S|$ as the ground state localizes onto few
configurations. The relationship between an orbital invariant and a basis variant cannot be otherwise.

\paragraph{An obstruction, not a weakness.} The failure is structural, not special to $\mathcal F_1$.
For the \emph{raw} determinant support (the unthresholded Slater rank), \emph{no} continuous functional
of a fixed set of reduced density matrices can lower-bound it (App.~\ref{app:theorem}): rank is a
coefficient-insensitive count, so $\ket{\Psi_\epsilon}=\ket{D_0}+\epsilon\sum_{k=1}^{K}\ket{D_k}$ has all
reduced density matrices converging to those of a single determinant---every continuous magic or
correlation functional $\to0$---while the rank stays $\Theta(K)$. This continuity route does \emph{not}
by itself carry over to the operative, coefficient-\emph{sensitive} $|\mathcal S|_\varepsilon$, which
collapses to $1$ as $\epsilon\to0$ in step with the functionals; the threshold-robust obstruction for
$|\mathcal S|_\varepsilon$ is supplied instead by the symmetry argument above---no
orbital-rotation-invariant functional can bound it---and by two \emph{equal-weight} witnesses, whose
amplitudes do not shrink. The two-determinant cat $(\ket{1\cdots10\cdots0}+\ket{0\cdots01\cdots1})/\sqrt2$
has $\gamma=\tfrac12\mathbb 1$, hence $\mathcal F_1=2n_{\rm o}$ (maximal) yet Slater rank $2$; and a paired (BCS)
Gaussian state has Gaussian rank $1$ yet fractional occupations. The only computable lower bounds on $|\mathcal S|$ are themselves
basis dependent---the orbital-Schmidt rank across a bipartition, i.e.\ the tensor-network bond
dimension~\cite{schollwock2011}, which we prove is a rigorous lower bound on the raw support,
$\chi\le|\mathcal S|$ (App.~\ref{app:theorem}, Prop.~\ref{prop:chibound}), and, consistently, the basis-dependent non-Gaussianity measures
(such as the natural-orbital participation entropy) that \emph{do} upper-bound the simulation
cost~\cite{tarabunga2026}---while a Gaussian-invariant magic measure is blind to them. We make the
obstruction constructive (App.~\ref{app:witness}, Thm.~\ref{thm:witness}): an antisymmetrized product of
$K$ strongly orthogonal geminals carries one-body magic $\mathcal F_1=2N_{\mathrm u}=4K$, higher-order
magic $\mathcal F_2=4K$, and pairing-basis determinant support $|\mathcal S|=2^{K}$ all growing without
bound, yet in that basis is a bond-dimension-$2$ tensor network---$O(1)$ bond dimension, $O(K)$
contraction and Born-sampling cost---so magic of \emph{every} order is unbounded at fixed $O(1)$ bond dimension. Because $\mathcal F_k$ is orbital-rotation invariant while the sampling
cost is not, no magic monotone is a sufficient certificate of hardness. Even the higher-order
antiflatness $\mathcal F_{k\ge2}$ fails: the witness has $\mathcal F_2=4K$ unbounded at $O(1)$ bond
dimension. What is hard is \emph{generic, high-rank} non-Gaussianity in the fixed operative basis---the
connected higher-body cumulants, which no orbital-invariant $\mathcal F_k$ can capture---and not the
paired stratum, whose amplitudes, overlaps, and number correlators are classically \emph{estimable} to
additive error under free-fermionic dynamics~\cite{oh2026}.
This bond dimension---which directly sets the classical cost of a matrix-product-state
simulation~\cite{schollwock2011}---is a rigorous lower bound on the \emph{raw} support,
$\chi\le|\mathcal S|$ (Prop.~\ref{prop:chibound}). We are explicit that the quantity we measure is the
operative thresholded support $|\mathcal S|_\varepsilon$, for which the relation is empirical rather than
proved: the support and the Schmidt rank are truncated at different tolerances ($99\%$ of the weight for
$|\mathcal S|$, $\varepsilon=0.05$ for $\chi$), so the thresholded support can fall below $\chi$, and it
does so in $7$ of the $38$ molecular points. What holds empirically is that, unlike the orbital-invariant
$\mathcal F_1$, the bond dimension co-varies with the measured cost: across both the nineteen-molecule
suite and open Hubbard chains, $\chi$ tracks $|\mathcal S|_\varepsilon$ (Spearman $\rho=0.90$ over the
$n=38$ molecular equilibrium-and-dissociation points and $\rho=0.57$ over the $n=30$ Hubbard points; both
$p<0.01$). The molecular set by itself does not discriminate between the two candidate resources---one-body
magic scores a comparable $\rho=0.86$ on the identical $n=38$ points, because there $\mathcal F_1$ and the
support both merely track static correlation. What discriminates is the Hubbard interaction sweep, where
$\chi$ holds at $\rho=0.57$ while $\mathcal F_1$ collapses to $\rho=-0.11$ ($n=30$, $p\approx0.56$, not
significant), increasing $U$ raising $\mathcal F_1$ while lowering $|\mathcal S|_\varepsilon$; that
contrast, not the molecular coefficient, is the evidence for the title. The absolute
$|\mathcal S|_\varepsilon$, sector dimension $D$, and $\chi$ underlying these coefficients are tabulated
per species in App.~\ref{app:repro}. Many-body
(bond-dimension) entanglement lower-bounds and, across the studied regimes, empirically tracks the
sampling cost; one-body entanglement---equivalently the one-body magic $\mathcal F_1$---does so only
where it coincides with static correlation. That single-particle entanglement can be maximal while
bipartite (bond) entanglement stays $O(1)$ is the fermionic counterpart of the magic--entanglement
decoupling seen in qubit matrix-product states~\cite{frautarabunga2024}.
Figure~\ref{fig:master} consolidates the picture over all $74$ exact ground states of this
work---nineteen molecules (equilibrium and dissociation), the Hubbard interaction sweep, and the
chain/ladder pair. The determinant support tracks the bond dimension in \emph{every} family (pooled
Spearman $\rho=0.72$, and positive within each), with the geminal witness ($\chi=2$, $|\mathcal S|=2^{K}$)
the deliberately loose lower bound; one-body magic does not, and fails in two visibly distinct ways---the
Hubbard branch runs the \emph{wrong} way ($U\!\uparrow$ raises $\mathcal F_1$ while lowering
$|\mathcal S|$) and the chain/ladder points sit at \emph{identical} $\mathcal F_1$ with $|\mathcal S|$
differing by up to $\sim\!2.2\times$. The cost is a property of the sampler's basis, read by $\chi$; the orbital
invariant $\mathcal F_1$ is blind to it.

\begin{figure}[htbp]\centering% ===== fig:master (REBUILT 2026-09-18) -- 74 exact ground states, native pgfplots.
% Emitted by p2_figs/make_figs.py from data/resource_master.json and p2_calc/stats_resource.json.
% No hand-typed numbers: every coordinate, coefficient and interval is read from those files.
\definecolor{rmMol}{HTML}{0072B2}\definecolor{rmHub}{HTML}{E69F00}
\definecolor{rmChn}{HTML}{009E73}\definecolor{rmLad}{HTML}{D55E00}\definecolor{rmInk}{HTML}{1B1B1B}
\definecolor{rmGrey}{HTML}{6E6E6E}
\begin{tikzpicture}
\begin{groupplot}[group style={group size=2 by 2, horizontal sep=1.8cm, vertical sep=2.05cm},
  width=7.45cm, height=5.6cm, paperaxis,
  title style={at={(0,1)},anchor=south west,font=\small\bfseries,yshift=2pt},
  label style={font=\normalsize}, tick label style={font=\footnotesize},
  legend style={font=\scriptsize,draw=black!25,fill=white,fill opacity=0.92,text opacity=1},
  legend cell align=left]
% ---- (a) |S| vs chi, with the eleven molecular floor points marked ----
\nextgroupplot[title={(a)~cost tracks $\chi$}, xmode=log, ymode=log,
  xlabel={minimal bond dimension $\chi$}, ylabel={determinant support $|\mathcal S|$},
  xmin=0.8,xmax=70,ymin=0.50,ymax=900000, ymajorgrids,xmajorgrids, grid style={draw=black!8},
  legend to name=fig2legend, legend columns=3, legend style={draw=black!25,/tikz/every even column/.append style={column sep=11pt}}]
\addplot[only marks,mark=*,mark size=1.9pt,rmMol,mark options={fill=rmMol,draw=white,line width=0.3pt}] coordinates {(2,2) (4,2) (4,4) (2,2) (2,5) (2,3) (14,8) (30,46) (7,7) (32,80) (11,13) (2,6) (8,6) (13,13) (42,68) (4,2) (8,11) (8,9) (36,91) (2,2) (4,5) (4,11) (2,4) (17,27) (7,9) (11,5) (7,25)};
\addlegendentry{molecules, $|\mathcal S|>1$ ($\times27$)}
\addplot[only marks,mark=o,mark size=2.2pt,rmMol,line width=0.8pt] coordinates {(1,1) (1,1) (1,1) (2,1) (1,1) (1,1) (1,1) (1,1) (4,1) (1,1) (1,1)};
\addlegendentry{molecules at the floor $|\mathcal S|=1$ ($\times11$)}
\addplot[only marks,mark=square*,mark size=1.9pt,rmHub,mark options={fill=rmHub,draw=white,line width=0.3pt}] coordinates {(12,330) (10,312) (6,260) (5,147) (4,68) (11,200) (10,196) (9,188) (9,168) (9,143) (11,3395) (11,3120) (8,2176) (8,880) (7,327) (16,2342) (16,2272) (15,2041) (14,1600) (13,1091) (17,35678) (14,30975) (10,17732) (6,5084) (6,1476) (13,26503) (13,25235) (13,21203) (11,13832) (9,7527)};
\addlegendentry{Hubbard sweep ($\times30$)}
\addplot[only marks,mark=triangle*,mark size=2.2pt,rmChn,mark options={fill=rmChn,draw=white,line width=0.3pt}] coordinates {(5,147) (8,880) (6,5084)};
\addlegendentry{chain ($\times3$)}
\addplot[only marks,mark=diamond*,mark size=2.2pt,rmLad,mark options={fill=rmLad,draw=white,line width=0.3pt}] coordinates {(13,224) (38,1479) (38,11375)};
\addlegendentry{2-leg ladder ($\times3$)}
\addplot[rmInk,densely dashed,line width=0.9pt,mark=none,forget plot] coordinates {(2,2) (2,4) (2,8) (2,16) (2,32)};
\node[rmInk,font=\scriptsize,anchor=west] at (axis cs:2.15,120000) {witness: $\chi{=}2,\;|\mathcal S|{=}2^{K}$};
\node[rmInk,font=\scriptsize,anchor=south east] at (rel axis cs:0.98,0.02) {pooled $\rho=+0.722$};
% ---- (b) |S| vs F1, same 74 points, same marking ----
\nextgroupplot[title={(b)~$\mathcal F_1$ does not}, ymode=log,
  xlabel={one-body magic $\mathcal F_1$}, ylabel={determinant support $|\mathcal S|$},
  xmin=-1.2,xmax=23,ymin=0.50,ymax=900000, ymajorgrids, grid style={draw=black!8}]
\addplot[only marks,mark=*,mark size=1.9pt,rmMol,mark options={fill=rmMol,draw=white,line width=0.3pt}] coordinates {(2.5114,2) (0.0750,2) (3.4603,4) (3.9849,2) (0.5181,5) (0.5731,3) (0.6882,8) (8.1924,46) (0.8474,7) (11.3765,80) (4.7428,13) (8.0009,6) (3.7542,6) (6.9358,13) (10.2890,68) (0.2515,2) (5.1131,11) (0.5214,9) (7.7351,91) (0.3864,2) (3.7026,5) (1.6728,11) (6.0001,4) (1.3066,27) (5.9378,9) (0.7900,5) (7.4487,25)};
\addplot[only marks,mark=o,mark size=2.2pt,rmMol,line width=0.8pt] coordinates {(0.0380,1) (0.0021,1) (0.0008,1) (0.0676,1) (0.0000,1) (0.0007,1) (0.0392,1) (0.0068,1) (0.0803,1) (0.0353,1) (0.0127,1)};
\addplot[only marks,mark=square*,mark size=1.9pt,rmHub,mark options={fill=rmHub,draw=white,line width=0.3pt}] coordinates {(0.6258,330) (2.2316,312) (5.9132,260) (9.6047,147) (11.2997,68) (0.3396,200) (1.1047,196) (2.8100,188) (4.9945,168) (6.5809,143) (0.8020,3395) (2.9184,3120) (7.8807,2176) (12.8163,880) (15.0697,327) (0.4855,2342) (1.6348,2272) (4.3683,2041) (7.9600,1600) (10.4394,1091) (0.9724,35678) (3.6002,30975) (9.8568,17732) (16.0306,5084) (18.8404,1476) (0.6329,26503) (2.1803,25235) (5.9969,21203) (11.0065,13832) (14.3023,7527)};
\addplot[only marks,mark=triangle*,mark size=2.2pt,rmChn,mark options={fill=rmChn,draw=white,line width=0.3pt}] coordinates {(9.6047,147) (12.8163,880) (16.0306,5084)};
\addplot[only marks,mark=diamond*,mark size=2.2pt,rmLad,mark options={fill=rmLad,draw=white,line width=0.3pt}] coordinates {(9.6948,224) (12.9107,1479) (16.0889,11375)};
\node[rmInk,font=\scriptsize,anchor=north west,align=left] at (rel axis cs:0.02,0.97) {pooled $\rho=+0.599$ over the same 74 rows\\Hubbard branch: $U\!\uparrow\;\Rightarrow\;\mathcal F_1\!\uparrow,\;|\mathcal S|\!\downarrow$};
% ---- (c) the four coefficients, the intervals the deposit supports, and the tie ceiling ----
\nextgroupplot[title={(c)~coefficients; intervals for $n{=}38$ only},
  xlabel={Spearman $\rho$ against $|\mathcal S|$}, ylabel={},
  xmin=-0.03,xmax=1.03,ymin=-1.55,ymax=5.45, xtick={0,0.25,0.5,0.75,1},
  ytick={1,2,3,4}, yticklabels={{$\mathcal F_1$, $n{=}74$},{$\chi$, $n{=}74$},{$\mathcal F_1$, $n{=}38$},{$\chi$, $n{=}38$}},
  y tick label style={font=\scriptsize}, ymajorgrids=false, xmajorgrids, grid style={draw=black!8}]
\addplot[rmGrey,densely dashed,line width=0.9pt,forget plot] coordinates {(0.9863,0.45) (0.9863,4.45)};
\addplot[rmMol,line width=1.0pt,forget plot] coordinates {(0.8051,4) (0.9589,4)};
\addplot[rmMol,line width=1.0pt,forget plot] coordinates {(0.8051,3.87) (0.8051,4.13)};
\addplot[rmMol,line width=1.0pt,forget plot] coordinates {(0.9589,3.87) (0.9589,4.13)};
\addplot[rmLad,line width=1.0pt,forget plot] coordinates {(0.7091,3) (0.9294,3)};
\addplot[rmLad,line width=1.0pt,forget plot] coordinates {(0.7091,2.87) (0.7091,3.13)};
\addplot[rmLad,line width=1.0pt,forget plot] coordinates {(0.9294,2.87) (0.9294,3.13)};
\addplot[only marks,mark=*,mark size=2.4pt,rmMol,mark options={fill=rmMol,draw=white,line width=0.4pt},forget plot] coordinates {(0.9011,4)};
\addplot[only marks,mark=*,mark size=2.4pt,rmLad,mark options={fill=rmLad,draw=white,line width=0.4pt},forget plot] coordinates {(0.8618,3)};
\addplot[only marks,mark=*,mark size=2.4pt,rmMol,mark options={fill=rmMol,draw=white,line width=0.4pt},forget plot] coordinates {(0.7224,2)};
\addplot[only marks,mark=*,mark size=2.4pt,rmLad,mark options={fill=rmLad,draw=white,line width=0.4pt},forget plot] coordinates {(0.5992,1)};
\node[rmGrey,font=\scriptsize,anchor=east] at (axis cs:0.9863,5.05) {tie ceiling $0.986$\;};
\node[rmInk,font=\scriptsize,anchor=west] at (axis cs:-0.01,-0.40) {$n{=}38$: $\Delta\rho=0.039$ ($t=0.94$, $p=0.35$)};
\node[rmInk,font=\scriptsize,anchor=west] at (axis cs:-0.01,-1.15) {$n{=}74$: $\Delta\rho=0.123$ ($t=1.50$, $p=0.14$)};
% ---- (d) the molecular coefficient in rank space: the eleven-point tied block ----
\nextgroupplot[title={(d)~the 38 molecules in rank space},
  xlabel={midrank of $|\mathcal S|$}, ylabel={midrank of the predictor},
  xmin=0,xmax=40,ymin=0,ymax=59, xmajorgrids, ymajorgrids, grid style={draw=black!8},
  ytick={0,10,20,30,40},
  legend style={at={(0.98,0.985)},anchor=north east,font=\scriptsize}]
\addplot[rmGrey,densely dashed,line width=0.8pt,forget plot] coordinates {(1,1) (38,38)};
\addplot[rmGrey,line width=0.7pt,forget plot] coordinates {(6.0,25.0) (6.0,41.5)};
\addplot[only marks,mark=*,mark size=1.9pt,rmMol,mark options={fill=rmMol,draw=white,line width=0.3pt}] coordinates {(6.0,5.0) (14.0,13.5) (6.0,5.0) (14.0,20.5) (6.0,5.0) (18.5,20.5) (6.0,13.5) (14.0,13.5) (6.0,5.0) (6.0,5.0) (21.0,13.5) (17.0,13.5) (26.0,33.0) (35.0,35.0) (25.0,25.0) (37.0,36.0) (31.5,30.5) (23.5,13.5) (6.0,5.0) (23.5,28.0) (6.0,5.0) (31.5,32.0) (6.0,20.5) (36.0,38.0) (14.0,20.5) (29.5,28.0) (27.5,28.0) (38.0,37.0) (6.0,5.0) (14.0,13.5) (6.0,5.0) (21.0,20.5) (29.5,20.5) (18.5,13.5) (34.0,34.0) (27.5,25.0) (21.0,30.5) (33.0,25.0)};
\addlegendentry{$\chi$, $\rho=0.901$}
\addplot[only marks,mark=square*,mark size=1.9pt,rmLad,mark options={fill=rmLad,draw=white,line width=0.3pt}] coordinates {(6.0,8.0) (14.0,23.0) (6.0,4.0) (14.0,11.0) (6.0,3.0) (18.5,24.0) (6.0,10.0) (14.0,27.0) (6.0,1.0) (6.0,2.0) (21.0,15.0) (17.0,17.0) (26.0,18.0) (35.0,36.0) (25.0,20.0) (37.0,38.0) (31.5,28.0) (23.5,35.0) (6.0,9.0) (23.5,26.0) (6.0,5.0) (31.5,32.0) (6.0,12.0) (36.0,37.0) (14.0,13.0) (29.5,29.0) (27.5,16.0) (38.0,34.0) (6.0,7.0) (14.0,14.0) (6.0,6.0) (21.0,25.0) (29.5,22.0) (18.5,31.0) (34.0,21.0) (27.5,30.0) (21.0,19.0) (33.0,33.0)};
\addlegendentry{$\mathcal F_1$, $\rho=0.862$}
\node[rmInk,font=\scriptsize,anchor=north west,align=left] at (rel axis cs:0.02,0.985) {11 of 38 tied at $|\mathcal S|=1$\\midrank $6.0$: $|\rho|\le 0.986$};
% ---- D13: the grey rule marks the block of tied midranks; it now carries a label ----
\node[rmGrey,font=\scriptsize,anchor=west] at (axis cs:7.1,33.0) {tied block};
\end{groupplot}
\node[anchor=north,yshift=-1.25cm] at (group c1r2.south east) {\ref*{fig2legend}};
\end{tikzpicture}

\caption{\textbf{The resource thesis over $74$ exact ground states.} Nineteen molecules at equilibrium and
dissociation, a Hubbard interaction sweep, and a matched chain/two-leg-ladder pair (all exact
diagonalization). \textbf{(a)}~The determinant support $|\mathcal S|$ tracks the minimal bond dimension
$\chi$ in every family (pooled Spearman $\rho=0.72$; the geminal witness $\chi{=}2,\,|\mathcal S|{=}2^{K}$
is the deliberately loose lower bound). \textbf{(b)}~One-body magic $\mathcal F_1$ does not (pooled
Spearman $\rho=0.60$ for $\mathcal F_1$ vs $|\mathcal S|$ over the same $74$ states): the Hubbard
branch runs the wrong way ($U\!\uparrow\Rightarrow\mathcal F_1\!\uparrow,|\mathcal S|\!\downarrow$) and the
chain/ladder points share the same $\mathcal F_1$ at $|\mathcal S|$ differing by up to $\sim\!2.2\times$. Cost is lower-bounded and tracked by
the basis-dependent $\chi$, not by the orbital-invariant $\mathcal F_1$.}\label{fig:master}\end{figure}

\paragraph{What one-body magic does, and where the hardness is.} None of this makes $\mathcal F_1$
uninformative. It is a faithful, efficiently measurable detector of \emph{static} (multireference)
correlation---generalizing to fermionic non-Gaussianity, and to nineteen FCI-verified molecules, the
qubit non-stabilizerness peak of molecular bonding~\cite{sarkis2025} and the reported proportionality
between qubit non-stabilizerness and Hartree--Fock weight (correlation) for weak-to-moderate
correlation---itself reported to break down beyond the Coulson--Fischer point~\cite{seibert2026}
(Table~\ref{tab:molecules}). What it is
\emph{not} is a predictor of the classical \emph{cost} of sampling, which---as the $\chi$--$|\mathcal S|$
tracking above shows---is lower-bounded and tracked by the basis-dependent entanglement, and by the fixed basis a sampler
cannot rotate, not by the orbital invariant $\mathcal F_1$. Cost and quantum hardness are thus distinct
axes: the classical cost of the simulation is governed by entanglement, the bond dimension $\chi$,
whereas a genuine quantum advantage---a regime beyond \emph{every} classical method---would require
\emph{generic, high-rank} non-Gaussianity in the operative basis, the connected two- and higher-body
cumulants beyond the covariance~\cite{coffman2025}, and specifically not the paired stratum, whose
free-fermionic estimation primitives remain classically tractable~\cite{oh2026}; the orbital-invariant
one-body magic $\mathcal F_1$ certifies neither. Fermionic magic thus supplies the missing link between the resource theory of non-Gaussianity
and sample-based electronic-structure computation---but the link is a \emph{demarcation}: what is so
easily measured (invariant one-body magic) diagnoses correlation, while the cost lives in the
entanglement and any genuine advantage one stratum above, in the generic higher-order sector.
